\documentclass[11pt, letterpaper]{article}

\usepackage[top=1.5cm, bottom=1.5cm, left=1.9cm, right=1.9cm]{geometry}
\usepackage[T1]{fontenc}
\usepackage[utf8]{inputenc}
\usepackage{mathptmx}
\usepackage{microtype}
\usepackage{titlesec}
\usepackage[dvipsnames]{xcolor}
\usepackage{enumitem}
\usepackage{array}
\usepackage{calc}
\usepackage[hidelinks]{hyperref}

\definecolor{darkred}{RGB}{139,0,0}
\definecolor{accent}{RGB}{30,55,120}

\hypersetup{colorlinks=true, urlcolor=accent, linkcolor=accent,
  pdftitle={Haoming Wang -- CV (ByteDance Seed Long-Range Video Generation)}, pdfauthor={Haoming Wang}}

\pagestyle{empty}
\setcounter{secnumdepth}{0}
\setlength{\parindent}{0pt}
\pagenumbering{gobble}
\raggedright

\titleformat{\section}{\scshape\bfseries\large}{}{0pt}{}[\vspace{1pt}\titlerule]
\titlespacing{\section}{0pt}{10pt}{4pt}

\newenvironment{datedentry}[1]{%
  \par\noindent
  \begin{minipage}[t]{2.8cm}\raggedright\small #1\end{minipage}%
  \hspace{0.3cm}%
  \begin{minipage}[t]{\dimexpr\linewidth-3.1cm\relax}\raggedright
}{\end{minipage}\par\medskip}

\newenvironment{publist}{
  \begin{enumerate}[leftmargin=1.7em, itemsep=4pt, parsep=0pt, topsep=2pt, label=\arabic*.]
}{\end{enumerate}}

\newcommand{\impactbullet}[1]{%
  \par\vspace{1pt}%
  \begin{itemize}[leftmargin=1.2em, topsep=0pt, itemsep=0pt, parsep=0pt, partopsep=0pt]
    \item #1
  \end{itemize}}

\newcommand{\me}{\underline{\textbf{Haoming Wang}}}
\newcommand{\meco}{\underline{\textbf{Haoming Wang}} (co-author)}
\newcommand{\meanon}{\underline{\textbf{Haoming Wang}}}

\begin{document}

\noindent
\begin{minipage}[t]{0.58\linewidth}
  \vspace{0pt}
  {\fontsize{24pt}{26pt}\selectfont \textbf{HAOMING WANG}}\\[2pt]
  {\scshape\normalsize Ph.D.\ Candidate}
\end{minipage}%
\begin{minipage}[t]{0.42\linewidth}
  \vspace{0pt}
  \raggedleft\small
  +1 412-909-7571 \\
  \href{mailto:haw200@pitt.edu}{haw200@pitt.edu} \\
  Pittsburgh, PA, USA \\[2pt]
  \href{https://haomingwang645.github.io/}{Homepage}\ $\diamond$\ \href{https://scholar.google.com.hk/citations?user=bf6LCjUAAAAJ}{Scholar}\ $\diamond$\ \href{https://www.linkedin.com/in/haoming-wang-40a6a7239/}{LinkedIn}\ $\diamond$\ \href{https://github.com/HaomingWang645}{GitHub}
\end{minipage}

\section{Research Interests}
Video Benchmark Synthesis \& Evaluation, Diffusion / Generative Image \& Video Models, Latent-Representation Shaping for VLMs, Cross-Frame Temporal Consistency.

\section{Education}

\begin{datedentry}{Sept 2022 -- May 2027 (Anticipated)}
\textbf{Ph.D.\ in Electrical and Computer Engineering}\\
University of Pittsburgh, Pittsburgh, PA, USA\\
Advisor: Prof.\ Wei Gao, Intelligent System Lab
\end{datedentry}

\begin{datedentry}{Sept 2018 -- May 2022}
\textbf{B.Eng.\ in Automation} \hfill GPA: 3.8/4.0\\
Zhejiang University, Hangzhou, China\\
Chu Kochen Honors College
\end{datedentry}

\section{Experience}

\begin{datedentry}{Sept 2022 -- Present}
\textbf{Graduate Student Researcher}\\
Intelligent System Lab, University of Pittsburgh, Pittsburgh, PA, USA
\begin{itemize}[leftmargin=1.1em, itemsep=2pt, topsep=3pt]
  \item \textbf{(May 2025 -- Present)} Led research on procedural / LLM-agentic 3D video benchmark synthesis (CVPR'26 Oral); Dirichlet-energy latent shaping of VLM 3D-topology subspaces; cross-frame consistency under sequential observation budgets.
  \item \textbf{(Nov 2024 -- May 2025)} Drove research on diffusion / generative-model interpretability and personalization: training-free fine-grained natural-language explanations of Stable Diffusion / GAN / autoregressive image generators (FineXL); explainable on-device LLM personalization (MobiSys'25 Xpert).
  \item \textbf{(Sept 2022 -- Oct 2024)} Designed federated learning systems for heterogeneous AIoT fleets: gradient correction under data / device heterogeneity (MobiCom'25); federated learning with unbounded device staleness (AAAI'25).
\end{itemize}
\end{datedentry}

\section{Publications (Peer-Reviewed)}

\begin{publist}
\item \textbf{\textcolor{darkred}{[CVPR'26 Oral]}\ InfiniBench: Infinite Benchmarking for Visual Spatial Reasoning with Customizable Scene Complexity}\\
\me, Qiyao Xue, Wei Gao.\\
\textit{IEEE/CVF Conference on Computer Vision and Pattern Recognition (CVPR)}, 2026 (Acceptance 25.4\%, Oral). \href{https://arxiv.org/abs/2511.18200}{[arXiv]}\,\href{https://github.com/pittisl/infinibench}{[Code]}\,\href{https://huggingface.co/datasets/Haoming645/infinibench}{[Data]}
\impactbullet{\textbf{Architected} a customizable framework that turns plain-language scene descriptions into photo-realistic 3D walk-through videos (Blender Cycles) with controllable identity and layout; established a new evaluation protocol for video-based reasoning -- a transferable methodology for long-video quality assessment.}

\item \textbf{\textcolor{darkred}{[MobiSys'25]}\ Never Start from Scratch: Expediting On-Device LLM Personalization via Explainable Model Selection}\\
\me, Boyuan Yang, Xiangyu Yin, Wei Gao.\\
\textit{Proceedings of the 23rd ACM International Conference on Mobile Systems, Applications and Services (MobiSys)}, 2025 (Acceptance 18.0\%). \href{https://haomingwang645.github.io/pdfs/xpert_paper.pdf}{[Paper]}
\impactbullet{\textbf{Developed Xpert}, an explainable model-selection method warm-starting on-device LLM personalization from the closest pre-personalized checkpoint.}

\item \textbf{\textcolor{darkred}{[MobiCom'25]}\ When Device Delays Meet Data Heterogeneity in Federated AIoT Applications}\\
\me, Wei Gao.\\
\textit{Proceedings of the 31st ACM International Conference on Mobile Computing and Networking (MobiCom)}, 2025 (Acceptance 17.1\%). \href{https://haomingwang645.github.io/pdfs/feddc_paper.pdf}{[Paper]}
\impactbullet{\textbf{Engineered} a federated AIoT training method that jointly disentangles and corrects for data heterogeneity and device delays via targeted gradient correction.}

\item \textbf{\textcolor{darkred}{[AAAI'25]}\ Tackling Intertwined Data and Device Heterogeneities in Federated Learning with Unlimited Staleness}\\
\me, Wei Gao.\\
\textit{Proceedings of the 39th AAAI Conference on Artificial Intelligence}, 2025 (Acceptance 23.4\%). \href{https://arxiv.org/abs/2309.13536}{[arXiv]}
\impactbullet{\textbf{Designed} a gradient-inversion technique that recovers usable training signal from arbitrarily stale device updates, enabling FL to converge on real heterogeneous fleets.}
\end{publist}

\section{Preprints \& Papers Under Review}

\begin{publist}
\item \textbf{Deciphering Personalization: Towards Fine-Grained Explainability in Natural Language for Personalized Image Generation Models (FineXL)}\\
\me, Wei Gao.\\
\textit{Under Review}, 2025. \href{https://haomingwang645.github.io/pdfs/finexl_paper.pdf}{[Paper]}
\impactbullet{\textbf{Architected FineXL}, a training-free method decomposing a personalized image generator's divergence from its base into orthogonal natural-language concepts; cut single-aspect explanation error by 56\% across \textbf{diffusion, GAN, and autoregressive} generators -- a new methodology for fine-grained perceptual evaluation of generative-model differences.}

\item \textbf{FreezeAsGuard: Mitigating Illegal Adaptation of Diffusion Models via Selective Tensor Freezing}\\
Kai Huang, \meco, Wei Gao.\\
\textit{Preprint}, 2024. \href{https://haomingwang645.github.io/pdfs/freezeasguard_paper.pdf}{[Paper]}\,\href{https://github.com/pittisl/FreezeAsGuard}{[Code]}
\impactbullet{\textbf{Co-developed} a bilevel-optimization defense that learns a tensor-freezing mask over Stable Diffusion -- a clean differentiable recipe for fine-grained manipulation of diffusion weights at scale; 37--47\% stronger mitigation than unlearning baselines.}

\item \textbf{Uncovering and Shaping the Latent Representation of 3D Scene Topology in Vision-Language Models}\\
\me, Wei Gao.\\
\textit{Preprint}, 2026. \href{https://haomingwang645.github.io/pdfs/arxiv26_vlm_latent_shaping.pdf}{[Paper]}\,\href{https://github.com/pittisl/vlm-latent-shaping}{[Code]}
\impactbullet{\textbf{Discovered} a VLM latent subspace encoding 3D scene topology and \textbf{introduced} a Dirichlet-energy regularizer beating SFT by up to 12.1\% -- a concrete instance of shaping hierarchical latent structure relevant to extended-generation latent design.}

\item \textbf{MosaicThinker: On-Device Visual Spatial Reasoning for Embodied AI via Iterative Construction of Space Representation}\\
\me, Qiyao Xue, Weichen Liu, Wei Gao.\\
\textit{Under Review}, 2026. \href{https://www.arxiv.org/abs/2602.07082}{[arXiv]}
\impactbullet{\textbf{Built} a training-free pipeline that stitches multi-view first-person video frames into a compact semantic map; replaced drift-prone sequential alignment with an MST routing scheme anchoring every frame to a global reference -- a transferable idea for long-horizon temporal consistency.}

\item \textbf{Reasoning Path and Latent State Analysis for Multi-view Visual Spatial Reasoning: A Cognitive Science Perspective}\\
Qiyao Xue, \meco, Weichen Liu, Shiqi Wang, Yuyang Wu, Wei Gao.\\
\textit{Under Review}, 2026. \href{https://arxiv.org/abs/2512.02340}{[arXiv]}
\impactbullet{\textbf{Co-authored} ReMindView-Bench (50K+ VQAs across 15 VLMs); demonstrated internal uncertainty cleanly separates correct vs.\ incorrect reasoning trajectories -- a candidate perceptual / coherence metric for assessing temporal consistency in generated video.}

\item \textbf{HEVEN \& GlobalNav (VLM-based autonomous mobile systems)}\\
Anonymous Authors (incl.\ \meanon, 3rd author).\\
\textit{Double-Blind Under Review}, 2026.
\impactbullet{\textbf{Co-authored} two VLM-based mobile systems: HEVEN (hierarchical spatial semantic tree + Tree-of-Thought planner; 73.86\% success on a real quadruped) and GlobalNav (infrastructure-assisted navigation, 90.91\% success in simulation).}

\item \textbf{Spatial Reasoning in Multimodal Large Language Models: A Survey of Tasks, Benchmarks and Methods}\\
Weichen Liu, Qiyao Xue, \me, Xiangyu Yin, Boyuan Yang, Wei Gao.\\
\textit{Under Review}, 2026.
\end{publist}

\section{Honors \& Awards}

\begin{datedentry}{April 2026}
\textbf{ECE Research Assistant of the Year}\\
Department of Electrical and Computer Engineering, University of Pittsburgh
\end{datedentry}

\section{Technical Skills}
\textbf{Generative Models:} Stable Diffusion, ControlGAN, Anole-7B, CLIP / FG-CLIP; bilevel mask optimization, linear-representation decomposition, Laplacian-eigenmap latent probing.\\[2pt]
\textbf{Video / 3D Pipelines:} Blender Cycles, Infinigen, Habitat, Isaac Sim; frontier-based camera planning, MatchAnything, DepthPro, SAM-3D / 3DGS.\\[2pt]
\textbf{Programming:} Python, PyTorch, CUDA, Hugging Face, ONNX, TensorRT.

\end{document}
